\documentclass{report}

\usepackage[a4paper, left=2cm, right=2cm, top=3cm, bottom=3cm]{geometry}
\usepackage{graphicx} % Required for inserting images
\usepackage{amsmath} %usa questo per scrivere le formule matematiche
\usepackage{parskip} %Questo pack elimina l'indentazione fastidiosa 
\usepackage{amsthm} % Questo pack permette di scrivere in modo carino i teoremi,
\usepackage{amssymb}
\usepackage{xcolor}    % Per i colori
\usepackage{framed}    % Per le cornici
\usepackage{mdframed}  % Cornici più avanzatedefinizioni, proposizioni ecc.

%Quelli che seguono sono i risultati dell'uso del pack amsthm
\newtheorem{teo}{Teorema}
\newtheorem{prop}{Proposizione}
\newtheorem{corollario}{Corollario}[teo] % Numerato in base ai teoremi
\newtheorem{defi}{Definizione}
\newtheorem{oss}{Osservazione}
\newtheorem{nota}{Nota}
\newtheorem{es}{Esempio}

% Definizione di un ambiente "Dimostrazione" senza numerazione
\renewenvironment{proof}{{\noindent\bfseries Dimostrazione. }}{\qed}

\newenvironment{nothing}{}{}







\title{Logica Matematica}
\author{Fabio Colletti}
\date{February 2025}

\begin{document}

\maketitle


\begin{nothing}
    \begin{center}
\huge\textbf{{premesse}}    
\end{center}


\[E=mc^2\]

Divideremo il corso in \textbf{2 parti:}

\begin{enumerate}

\item Logica Proposizionale
\item Logica Predicativa

\end{enumerate}
    La Logica proposizionale si occupa di valutare la verità o la falsità di un insieme di elementi, senza andare a verificare cosa sono, quando e perché. 

La logica predicativa \emph{estende naturalmente} quella proposizionale, aggiungendo più contesto alle nostre teorie. 
\vspace{3cm}
\begin{center}
Immergiamoci in questo mare di divertimento!
\end{center} 
\end{nothing}

\newpage

\chapter{Logica Proposizionale}
\begin{nothing}
    Iniziamo con un semplice esempio.
Immagina di aver appena visto l'uomo dei tuoi sogni, alto, moro, riccio, e, come direbbe Vannacci, di etnia non italica.

Cerchiamo ora di \emph{formalizzare} queste informazioni.
Per farlo utilizziamo 4 proposizioni:

    $P1=$ "L'uomo è alto"

    $P2=$ "L'uomo è moro"

    $P3=$ "L'uomo è riccio"

    $P4=$ "L'uomo non è di etnia italica"

In questo scenario, queste sono tutte affermazioni VERE.

Ora se inserissimo un'altra proprietà, come ad esempio:

$P5=$ "L'uomo è basso"

Otterremmo che questa proposizione è FALSA.

In generale, data una proposizione $P$, se $P$ è vera allora avrà valore 1, mentre se è falsa avrà valore 0. 
Formuleremo meglio questo concetto in seguito.

\bigskip
Bene, quello che faremo in logica proposizionale sarà fare questo (e qualcosa di più).
\end{nothing}
\section{Introduzione alla Logica}
\begin{defi}[Linguaggio Proposizionale]
    Un \textbf{Linguaggio Proposizionale} è un insieme di variabili proposizioni \textbf{L}=\{$P1, P2, P3,$ ...\}.

    Un linguaggio proposizionale può essere:
    \begin{itemize}
        \item Finito
        \item Numerabile
        \item Non numerabile
    \end{itemize}
    
\end{defi} %Linguaggio Proposizionale
\bigskip

\begin{nota}
    Una variabile proposizionale è una proposizione che non è ulteriormente scomponibile. \\E' come i numeri, ossia è un concetto primitivo, non ulteriormente definibile.
\end{nota} 
\bigskip

\begin{defi}[Connettivi Logici]
    I \textbf{Connettivi Logici} servono, appunto, a delle proposizioni, per formarne altre di più complesse.
    Questi sono: $\rightarrow$ , $\land$ , $\lor$ , $\neg$ .
    Il loro significato intuitivo è:
    \begin{itemize}
        \item $A\land B$ significa A e B
        \item $A\lor B$ significa A o B
        \item $\neg A$ significa non A
        \item $A\rightarrow B$ significa se A allora B
    \end{itemize}
\end{defi} %Connettivi Logici
\bigskip

\begin{nothing}
    Quello che vorremmo fare adesso è trovare un modo per dare dei valori di verità e di falsità alle nostre variabili.  \\Diamo per assodato che una proposizione o è vera o è falsa, non può essere entrambe e deve essere almeno una delle due.
Potrebbe sembrare ambiguo, ma tranquilli, la situazione peggiorerà.
\end{nothing}
\bigskip

\begin{defi}[Assegnamento]
    Un \textbf{Assegnamento} è una funzione $\alpha$ che associa ad ogni variabile proposizionale 0 oppure 1.
    \[\alpha : L\rightarrow\{0,1\}\ \]
    
\end{defi}
\bigskip %assegnamento

\begin{defi}[Insieme degli enunciati]
    Dato un linguaggio L, l'\textbf{Insieme delle Proposizioni} (o degli enunciati) del linguaggio è il più piccolo insieme chiuso sotto l'utilizzo dei connettivi logici.
\end{defi}
\bigskip

\begin{oss}[costruzione di P]
    Costruiamo per ricorsione l'insieme delle proposizioni. 
    \\Sia $L$ un linguaggio,Poniamo: 
    \\[\baselineskip]
    $P_0=$L
    \\$P_{n+1}=\{\neg A,\quad A\land B,\quad A\rightarrow B,\quad A\lor B\ |\ A,B\in P_n\}$
    \\$P=\bigcup\limits_{n=0}^{\infty} P_n $
    \\[\baselineskip]
    La verifica che $P$ sia chiuso rispetto ai connettivi logici è lasciata alla lettrice \\(se sei un lettORE non mi interessa che tu abbia questo file).
\end{oss}
\bigskip 

\begin{prop}[Estensione dell'assegnamento]
    Sia L un linguaggio e $\alpha$ un suo assegnamento. Allora possiamo definire un assegnamento $\tilde\alpha : P \rightarrow \{0,1\}$ coerente.
\end{prop}
\begin{proof}
Sia $\alpha$ un asegnamento di $L$ e $F\in P$, definiamo l'assegnamento per induzione:
\begin{itemize}
    \item Se $F = p$ per qualche $p\in L$, allora è già definito.
    \item Se $F=\neg G$ definiamo $\tilde\alpha(\neg G)=$ 
    $\begin{cases}
        0, & \text{se } \alpha(G)=1 \\ 
        1, & \text{se } \alpha(G)=0
    \end{cases}$
    \item così via per gli altri 3 connettivi logici.
\end{itemize}
\end{proof}
\bigskip

\begin{oss}
    Sia L un linguaggio, P il suo insieme delle proposizioni e $F\in P$, $p_1,..,p_n$ le variabili proposizionali che compaiono in F. Se $\alpha$ e $\beta$ sono due assegnamenti di L che concordano du quelle variabili, allora $\tilde\alpha(F) = \tilde\beta(F)$.
\end{oss}
\bigskip

\begin{nota}
    Quando non è necessario specificarlo, faremo abuso di notazione e chiamaremo allo stesso modo l'assegnamento sulle variabili proposizionali e quello sulle proposizioni.\\
    Del resto faremo abuso di notazione e "confonderemo" un linguaggio con il suo insieme delle proposizioni.\bigskip

    N.B: Una Proposizione è SEMPRE finita!
\end{nota}
\bigskip

\begin{defi}[Forma Normale]
    E' facile convincersi (al massimo con qualche esempio), che è possibile "esprimersi" usando 3 dei 4 connettivi definiti prima, scartando quello dell'implicazione. \\
    Usare SOLO $\neg,\ \lor,\ \land$ è detta \textbf{Forma Normale}
\end{defi}
\bigskip

\begin{teo}[Completezza Funzionale]
    Sia L un linguaggio numerabile, sia $f:\{0,1\}^n \rightarrow \{0,1\}$. \\
    Esiste una proposizione A contenente n var proposizionali $p_1,...,p_n$ t.c. $\alpha(A)=f(\alpha(p_1),...,\alpha(p_n)),\ \forall\alpha$ assegnamento.
\end{teo}
\begin{proof}
    dimostriamo per induzione.
    \begin{itemize}
        \item Base induttiva: n=1.
        \item Supponiamo vero per n.
        \item Dimostriamo per n+1.
    \end{itemize}
\end{proof}
\bigskip

\section{Implicazione Logica \& Soddisfacibilità}
\begin{defi}[Implicazione Logica]
    Siano $A_1,...,A_n,A_{n+1}$ proposizioni. \\Diciamo che $A_1,...,A_n$ \textbf{Implicano Logicamente} $A_{n+1}$ sse:
    
    $\forall \alpha\ | \ \alpha(A_1)=...=\alpha(A_n)=1 \,$ si ha che $\alpha(A_{n+1})=1 $ 

    In tal caso scriveremo: $A_1,...,A_n \models A_{n+1}$
\end{defi}
\bigskip

\begin{defi}.\par
   \begin{itemize}
       \item A sat sse $\exists \alpha \ | \ \alpha(A)=1 $
       \item A unsat sse $ \nexists \alpha \ | \ \alpha(A)=1$
       \item A tatut sse $\forall\alpha, \ \alpha(A)=1$ 
   \end{itemize} 
\end{defi}
\bigskip

\begin{defi}[Teoria]
    Sia L un lingaggio, una \textbf{Teoria} è un insieme di proposizioni T=\{A1, A2, ...\}. T può essere:
    \begin{itemize}
        \item Finita
        \item Numerabile
        \item Non numerabile
    \end{itemize}
\end{defi}
\bigskip

\begin{prop}
    Sono equivalenti:
    \begin{enumerate}
        \item T è sat
        \item $\forall A, se \ T\models A \ \rightarrow T \not\models A$
        \item $\exists A \ t.c. \ T\not\models A$
    \end{enumerate}
\end{prop}

\begin{proof}. \par
(1 $\Rightarrow$ 2) \\
Sia T è soddisfacile Sia A t.c. $T \models A$ \\
$\Rightarrow \ \forall a \ t.c. \ a(T)=1, \ a(A)=1$ \\
$\Rightarrow \ \forall a \ t.c. \ a(T)=1, \ a(\neg A)=0$
$\Rightarrow \ T \not\models A$

(2 $\Rightarrow$ 3) \\
Sia $A \ t.c. \ T \models A$ \\
$\Rightarrow T \not\models \neg A$

(3 $\Rightarrow 1$) \\
dimostriamo con la contronominale.
Se T è unsat allora segue qualunque cosa, dunque T implica logicamente ogni proposizione.
\end{proof}
\bigskip






\section{Compattezza \& Conseguenze}

\begin{teo}[Erdős--de Bruijn]
Sia $G = (V, E)$ un grafo (eventualmente infinito). Se ogni sottografo finito di $G$ è $k$-colorabile, allora $G$ è $k$-colorabile.
\end{teo}

\begin{proof}
Per ogni vertice $v \in V$ e ogni colore $i \in \{1, \dots, k\}$, introduciamo una proposizionale $P_{v,i}$, da interpretare come ``il vertice $v$ ha colore $i$''.

Definiamo una teoria proposizionale $\Gamma$ contenente le seguenti formule:

\begin{itemize}
  \item[\textbf{(A)}] Assegnamento di almeno un colore:
  \[
  \bigvee_{i=1}^k P_{v,i} \quad \text{per ogni } v \in V.
  \]
  
  \item[\textbf{(U)}] Unicità del colore:
  \[
  \lnot (P_{v,i} \land P_{v,j}) \quad \text{per ogni } v \in V,\; i \ne j.
  \]
  
  \item[\textbf{(D)}] Diversità tra vertici adiacenti:
  \[
  \lnot (P_{u,i} \land P_{v,i}) \quad \text{per ogni } \{u,v\} \in E,\; i = 1, \dots, k.
  \]
\end{itemize}

Sia $\Gamma_0 \subseteq \Gamma$ un sottoinsieme finito. Allora $\Gamma_0$ coinvolge solo un numero finito di vertici e archi, e quindi corrisponde a un sottografo finito $G_0 \subseteq G$.

Per ipotesi, ogni sottografo finito di $G$ è $k$-colorabile. Dunque $\Gamma_0$ è soddisfacibile.

Per il Teorema di Compattezza, $\Gamma$ è soddisfacibile.

Pertanto esiste un'assegnazione dei colori ai vertici di $G$ che soddisfa $\Gamma$, ossia una $k$-colorazione valida di $G$.
\end{proof}



\bigskip
\bigskip
\bigskip
    
\begin{teo}[Teorema delle Estensioni Totali]
Sia $(X, <)$ un ordine parziale. Allora esiste un ordine totale $\prec$ su $X$ tale che
\[
\forall x, y \in X,\; x < y \Rightarrow x \prec y.
\]
\end{teo}

\begin{proof}
Introduciamo, per ogni coppia distinta $x, y \in X$, una proposizionale $P_{x,y}$, da interpretare come ``$x \prec y$''.

Definiamo una teoria proposizionale $\Gamma$ composta dalle seguenti formule:

\begin{itemize}
  \item[\textbf{(T)}] Totalità:
  \[
  P_{x,y} \lor P_{y,x} \quad \text{per ogni } x \ne y \in X.
  \]
  
  \item[\textbf{(A)}] Antisimetricità:
  \[
  \lnot (P_{x,y} \land P_{y,x}) \quad \text{per ogni } x \ne y \in X.
  \]
  
  \item[\textbf{(Tr)}] Transitività:
  \[
  (P_{x,y} \land P_{y,z}) \rightarrow P_{x,z} \quad \text{per ogni } x, y, z \in X \text{ distinti}.
  \]
  
  \item[\textbf{(E)}] Estensione dell’ordine parziale:
  \[
  P_{x,y} \quad \text{per ogni } x < y.
  \]
\end{itemize}

Sia $\Gamma_0 \subseteq \Gamma$ un sottoinsieme finito. Allora $\Gamma_0$ coinvolge solo un insieme finito $X_0 \subseteq X$. L'ordine parziale $<$ ristretto a $X_0$ può essere esteso a un ordine totale $\prec_0$ su $X_0$, che soddisfa tutte le formule di $\Gamma_0$.

Infatti:


Quindi ogni sottoinsieme finito $\Gamma_0$ è soddisfacibile.

Per il Teorema di Compattezza, anche $\Gamma$ è soddisfacibile. Pertanto, esiste un ordine totale $\prec$ su $X$ che soddisfa $\Gamma$, e in particolare estende l'ordine parziale $<$.
\end{proof}

\bigskip
\bigskip
\bigskip


\begin{teo}[Lemma di Kőnig (via compattezza)]
Sia $T \subseteq \{0,1\}^{<\omega}$ un albero binario infinito, ovvero:

\begin{itemize}
  \item $\varepsilon \in T$ (dove $\varepsilon$ è la stringa vuota),
  \item se $s \in T$ e $t \preceq s$ (prefisso), allora $t \in T$,
  \item per ogni $n \in \mathbb{N}$, l'insieme $\{s \in T : |s| = n\}$ è finito e non vuoto.
\end{itemize}

Allora esiste una funzione $f : \mathbb{N} \to \{0,1\}$ tale che per ogni $n$, $f \upharpoonright n \in T$.
\end{teo}

\begin{proof}
Per ogni $n \in \mathbb{N}$, siano $T_n = \{ s \in T : |s| = n \}$ i nodi al livello $n$.

Introduciamo proposizionali $P_s$ per ogni $s \in T$, da interpretare come ``il cammino scelto passa per $s$''.

Definiamo una teoria proposizionale $\Gamma$ contenente le seguenti formule:

\begin{itemize}
  \item[\textbf{(L)}] Ad ogni livello $n$ si sceglie esattamente un nodo:
  \[
  \bigvee_{s \in T_n} P_s \quad \text{e} \quad \lnot(P_s \land P_t) \quad \text{per } s, t \in T_n,\; s \ne t.
  \]
  
  \item[\textbf{(C)}] Coerenza per prefissi:
  \[
  P_s \rightarrow P_t \quad \text{per ogni } s, t \in T,\; t \preceq s.
  \]
\end{itemize}

Sia $\Gamma_0 \subseteq \Gamma$ un sottoinsieme finito. Esiste un massimo livello $n$ tale che $\Gamma_0$ contiene solo proposizionali $P_s$ con $|s| \leq n$. Poiché $T_n$ è finito e non vuoto, possiamo scegliere un cammino finito \( s_0 \prec s_1 \prec \dots \prec s_n \) in \( T \), e assegnare verità a \( P_{s_i} \) coerentemente con $\Gamma_0$.

Quindi ogni sottoinsieme finito $\Gamma_0$ è soddisfacibile. Per il Teorema di Compattezza, anche $\Gamma$ è soddisfacibile.

Una valutazione che soddisfa $\Gamma$ seleziona per ogni $n$ esattamente un nodo $s_n \in T_n$, con $s_n \preceq s_{n+1}$. Definendo $f(n) = s_{n+1}(n)$ otteniamo una funzione $f : \mathbb{N} \to \{0,1\}$ tale che ogni prefisso $f \upharpoonright n \in T$.
\end{proof}



\section{Completezza}

\begin{prop}[Correttezza]
    tutti gli assiomi sono tautologie
\end{prop}

\begin{proof}
    Ogni istanza di un assioma `e una tautologia, e il Modus Ponens conserva la verità: se A è vera e A → B `e vera, allora B è vera.    
\end{proof}

\bigskip


\begin{teo}[Completezza della logica proposizionale]
Sia $\Gamma$ una teoria proposizionale (insieme di formule proposizionali), e sia $\varphi$ una formula proposizionale tale che:
\[
\Gamma \models \varphi
\]
cioè ogni valutazione che soddisfa tutte le formule di $\Gamma$ soddisfa anche $\varphi$.

Allora:
\[
\Gamma \vdash \varphi.
\]
\end{teo}

\begin{proof}
Supponiamo per assurdo che $\Gamma \not\vdash \varphi$.

Allora l'insieme $\Gamma \cup \{\lnot\varphi\}$ è **consistente**, ovvero **non** dimostra una contraddizione:
\[
\Gamma \cup \{\lnot\varphi\} \not\vdash \bot.
\]

Poiché la logica proposizionale è finitamente consistente (ogni dimostrazione è finita), ciò equivale a dire che **ogni sottoinsieme finito** di $\Gamma \cup \{\lnot\varphi\}$ è consistente.

Per il Teorema di Completezza finita (o più precisamente: per il Teorema di Compattezza proposizionale), da tale consistenza segue che:
\[
\Gamma \cup \{\lnot\varphi\} \text{ è soddisfacibile}.
\]

Cioè esiste una valutazione booleana $v$ tale che:
\[
v(\psi) = 1 \quad \text{per ogni } \psi \in \Gamma,\quad \text{e} \quad v(\varphi) = 0.
\]

Ma ciò contraddice l’ipotesi iniziale $\Gamma \models \varphi$, secondo cui ogni valutazione che soddisfa $\Gamma$ deve soddisfare anche $\varphi$.

Segue che $\Gamma \vdash \varphi$.
\end{proof}





\chapter{Logica Predicativa}
\section{Introduzione alla Logica Predicativa}




\end{document}

